\documentclass[11pt]{article}
\usepackage[T1]{fontenc}
\usepackage{lmodern}
\usepackage[a4paper,margin=28mm]{geometry}
\usepackage{amsmath,amssymb,amsthm}
\usepackage{microtype}
\usepackage[hidelinks]{hyperref}
\hypersetup{pdftitle={Sharp connection-Dirac spectral-sum domination for graph complexes},pdfauthor={Xiaojian Zeng}}
\newtheorem{theorem}{Theorem}
\newtheorem{lemma}{Lemma}
\newcommand{\tr}{\operatorname{tr}}
\newcommand{\rank}{\operatorname{rank}}
\newcommand{\spec}{\operatorname{spec}}
\newcommand{\R}{\mathbb{R}}
\newcommand{\T}{\mathsf{T}}
\title{Sharp connection--Dirac spectral-sum domination\\for graph complexes}
\author{Xiaojian Zeng\\[2pt]
\small PBC School of Finance, Tsinghua University, 100190\\
\small \href{mailto:xiaojianzeng@foxmail.com}{xiaojianzeng@foxmail.com}}
\date{}
\begin{document}
\maketitle

\begin{abstract}
Let $L$ be the face-intersection matrix of a nonempty finite simplicial
complex of dimension at most one, and let $D$ be its signed incidence Dirac
matrix. We prove that the sum of the $k$ largest eigenvalues of $L$ exceeds
the corresponding sum for $D$ by at least one, for every admissible $k$.
The constant is sharp for a single edge. The proof combines the unsigned
incidence spectral transform with degree majorization, a cycle estimate,
and a trace--rank bound. Disconnected graphs, isolated vertices and
non-bipartite graphs are included. Knill's original comparison in arbitrary
dimension remains unresolved.
\end{abstract}

\section{The graph-complex case}

For a real symmetric $n\times n$ matrix $M$, write its eigenvalues, with
multiplicity, as $\lambda_1(M)\geq\cdots\geq\lambda_n(M)$ and set
\[
 S_k(M)=\sum_{i=1}^k\lambda_i(M),\qquad 1\leq k\leq n.
\]
Let $G$ be a finite abstract simplicial complex: a family of nonempty sets
closed under taking nonempty subsets. Its connection matrix and Dirac matrix
are
\[
 (L_G)_{xy}=\mathbf{1}_{x\cap y\ne\varnothing},\qquad
 D_G=d+d^{\T},\qquad x,y\in G,
\]
where $d_{xy}$ is the oriented incidence sign if $y\subset x$ has codimension
one, and is zero otherwise. More explicitly, for an ordered face
$x=(v_0<\cdots<v_p)$ with $p\geq1$, the coefficient at
$y=x\setminus\{v_j\}$ is $(-1)^j$. The empty face is not a matrix index.
These are the conventions in Knill's comparison
\cite[Sections 1.1 and 5.1--5.3]{Knill2026}. Reorienting faces conjugates
$D_G$ by a diagonal sign matrix, so the spectral sums are orientation
independent.

Knill conjectures $S_k(L_G)\geq S_k(D_G)$ for arbitrary finite complexes
\cite[Conjecture 1]{Knill2026}. His earlier August 2025 post already proposes
the stronger additive-one inequality and its sharpness
\cite[August 17 update]{Knill2025}. We prove that proposed strengthening for
the entire graph-complex class. The original arbitrary-dimensional problem
remains unresolved.

\begin{theorem}\label{thm:main}
Let $G$ be a nonempty finite abstract simplicial complex whose faces all
have cardinality at most two. Then
\begin{equation}\label{eq:main}
 S_k(L_G)\geq S_k(D_G)+1\qquad(1\leq k\leq |G|).
\end{equation}
The constant $1$ cannot be increased: equality holds at $k=1$ for the
complex consisting of one edge and its two vertices.
\end{theorem}

The theorem allows any finite simple graph, with its vertices and edges as
faces; connectedness and bipartiteness are not assumed. In particular, a
triangle is allowed without inserting a filled two-dimensional face.

The closest spectral input is Knill's hydrogen identity
$L-L^{-1}=|D|_{\mathrm{ent}}^2$ for graph complexes, where
$|D|_{\mathrm{ent}}$ denotes entrywise absolute value
\cite[Theorem 1, Section 2.4]{Knill2018}. It concerns unsigned incidence;
the signed and unsigned Laplacians are conjugate in the bipartite case
\cite[Section 1.4]{Knill2018}. We derive the required spectral transform
directly below. To handle all partial sums for non-bipartite graphs, we
combine it with a degree-majorization estimate. This is the step beyond
the unsigned spectral information and the spectral-radius comparison in
\cite[Corollary 1, Section 2.7]{Knill2018}.

\section{Incidence spectra}

Write $v,e,c$ for the numbers of vertices, edges and connected components,
counting isolated vertices as components, and put $n=v+e$. Thus
$v\geq c\geq1$. Write $E$ for the edge set and abbreviate $L_G,D_G$ to
$L,D$. Let $B,C\in\R^{v\times e}$ be the unsigned and signed
vertex-by-edge incidence matrices. Each column of $B$ has two entries $1$;
the corresponding column of $C$ has one entry $-1$ and one entry $1$.
If $A$ is the adjacency matrix and $\Delta$ is the diagonal degree matrix,
then
\begin{equation}\label{eq:laplacians}
 Q=BB^{\T}=\Delta+A,\qquad H=CC^{\T}=\Delta-A.
\end{equation}
Let $q_1\geq\cdots\geq q_v\geq0$ and
$h_1\geq\cdots\geq h_v\geq0$ be their eigenvalues, and let
$\delta_1\geq\cdots\geq\delta_v\geq0$ be the vertex degrees in decreasing
order. With vertices preceding edges, the definitions give
\begin{equation}\label{eq:blocks}
 L=\begin{pmatrix}I_v&B\\ B^{\T}&B^{\T}B-I_e\end{pmatrix},\qquad
 D=\begin{pmatrix}0&C\\ C^{\T}&0\end{pmatrix}.
\end{equation}
Indeed, the diagonal entries of $B^{\T}B$ equal two, and distinct edges
share at most one vertex.

\begin{lemma}\label{lem:spectra}
Set
\[
 f(t)=\frac{t+\sqrt{t^2+4}}2\qquad(t\geq0).
\]
The matrix $L$ has exactly $v$ positive and $e$ negative eigenvalues, and
its positive eigenvalues are $f(q_1),\ldots,f(q_v)$ in decreasing order.
Consequently,
\begin{align}
 S_k(L)&=\sum_{i=1}^k f(q_i),&
 S_k(D)&=\sum_{i=1}^k\sqrt{h_i} &&(1\leq k\leq v),\label{eq:initial}\\
 S_k(L)&\geq n &&&&(v\leq k\leq n).\label{eq:trace}
\end{align}
\end{lemma}

\begin{proof}
For each positive singular value $s$ of $B$, a singular-value decomposition
reduces the corresponding part of $L$ to
\[
 \begin{pmatrix}1&s\\s&s^2-1\end{pmatrix}.
\]
Its characteristic polynomial is $z^2-s^2z-1$, so its eigenvalues are
$f(s^2)$ and $-1/f(s^2)$. The spaces $\ker B^{\T}$ and $\ker B$ give the
remaining eigenvalues $1$ and $-1$, respectively. Since $f$ is increasing
and $f(0)=1$, this proves the asserted positive list and inertia. For
$k\geq v$, all omitted eigenvalues are negative, and hence
$S_k(L)\geq\tr L=n$.

Similarly, the nonzero eigenvalues of $D$ are the pairs
$\pm\sqrt{h_i}$ for $h_i>0$. If $r=\rank C$, then there are $n-2r$
zeros. Because $r\leq e$, the number $n-r$ of nonnegative Dirac
eigenvalues is at least $v$. Thus the first $v$ positions use precisely
the positive square roots followed by enough zeros, as in
\eqref{eq:initial}. These arguments also apply when $e=0$.
\end{proof}

\section{Proof of Theorem \ref{thm:main}}

We first record two ordered comparisons:
\begin{equation}\label{eq:degree}
 \sum_{i=1}^kq_i\geq\sum_{i=1}^k\delta_i\quad(1\leq k\leq v),
 \qquad h_i\leq2\delta_i\quad(1\leq i\leq v).
\end{equation}
For the first, use the Ky Fan variational formula \cite{Fan1949},
\[
 S_k(M)=\max\{\tr(PM):P=P^{\T}=P^2,\ \rank P=k\},
\]
with the coordinate projection onto the $k$ largest diagonal entries of
$Q$. To recall its elementary proof, diagonalize $M$. In that basis the
diagonal entries of $P$ lie in $[0,1]$ and sum to $k$, so their weighted
eigenvalue sum is at most the sum of the largest $k$ eigenvalues; projection
onto the corresponding eigenspaces attains the bound. For the second
comparison, $2\Delta-H=Q\succeq0$ by \eqref{eq:laplacians}.
Applying the eigenvalue min--max principle to this quadratic-form inequality
gives $h_i\leq\lambda_i(2\Delta)=2\delta_i$.

Define
\[
 R(t)=\frac{\sqrt{t^2+4}-2}{2}\geq0,
 \qquad f(t)=1+\frac t2+R(t).
\]
For $1\leq k\leq v$, Lemma~\ref{lem:spectra} and \eqref{eq:degree} imply
\begin{align}
 S_k(L)-S_k(D)
 &\geq\sum_{i=1}^kR(q_i)+k+\frac12\sum_{i=1}^kq_i
                  -\sum_{i=1}^k\sqrt{2\delta_i}\notag\\
 &\geq R(q_1)+\frac12\sum_{i=1}^k
                    (\sqrt{\delta_i}-\sqrt2)^2
 \geq R(q_1).\label{eq:bridge}
\end{align}
The second line uses degree majorization and discards only nonnegative
terms $R(q_i)$.

\paragraph{Non-bipartite graphs, $1\leq k\leq v$.}
Such a graph contains an odd cycle. For example, coloring by distance
parity in each component gives either a bipartition or an odd closed walk,
from which an odd cycle can be extracted. Let its length be $\ell$, and
let $x$ be $1/\sqrt\ell$ on its vertices and zero elsewhere. Then
$\|x\|=1$, and the cycle edges alone give
\[
 q_1\geq x^{\T}Qx
      =\sum_{ab\in E}(x_a+x_b)^2\geq4.
\]
Since $R$ is increasing, \eqref{eq:bridge} gives
\[
 S_k(L)-S_k(D)\geq R(4)=\sqrt5-1>1
 \qquad(1\leq k\leq v).
\]

\paragraph{Bipartite graphs, $2\leq k\leq v$.}
Let $J$ be diagonal with entries $1$ and $-1$ on the two vertex classes;
choose either sign on each isolated vertex. Then $JHJ=Q$, so $h_i=q_i$,
including zero multiplicities. For every $t\geq0$,
\[
 f(t)\geq1+t/2\geq\sqrt t+1/2,
\]
because the second difference is $(\sqrt t-1)^2/2$. Therefore
\[
 S_k(L)-S_k(D)\geq k/2\geq1\qquad(2\leq k\leq v).
\]

\paragraph{The endpoint $k=1$.}
For every complex, the definitions give the entrywise inequality
$|D|_{\mathrm{ent}}+I\leq L$. For a real unit vector $u$, put $w=|u|$
coordinatewise. Then
\[
 u^{\T}Du\leq w^{\T}|D|_{\mathrm{ent}}w
 \leq w^{\T}(L-I)w\leq\lambda_1(L)-1.
\]
Maximizing over $u$ proves the endpoint. This also handles $v=1$.

\paragraph{The tail $v\leq k\leq n$.}
The identity
\[
 x^{\T}Hx=\sum_{ab\in E}(x_a-x_b)^2
\]
shows that $\ker H$ consists exactly of functions constant on each
connected component. Thus $r=\rank C=v-c$. Every partial sum of $D$ is
at most the sum of its positive eigenvalues. Using
$\sqrt t\leq(t+1)/2$ and $\tr H=2e$ gives
\[
 S_k(D)\leq\sum_{h_i>0}\sqrt{h_i}
 \leq\frac{2e+r}{2}=e+\frac{v-c}{2}.
\]
By \eqref{eq:trace},
\[
 S_k(L)-S_k(D)\geq n-e-\frac{v-c}{2}
                  =\frac{v+c}{2}\geq1\qquad(v\leq k\leq n).
\]
Together the ranges cover every $k$ in \eqref{eq:main}. For a nonempty
edgeless complex one can also check directly that $L=I_v$ and $D=0$.

\paragraph{Sharpness.}
For one edge and its two vertices, direct diagonalization gives
\[
 \spec(L)=(1+\sqrt2,1,1-\sqrt2),\qquad
 \spec(D)=(\sqrt2,0,-\sqrt2).
\]
The gap at $k=1$ is exactly one, completing the proof.\qed

\section{Computational checks and scope}

The proof above is independent of computation. As implementation checks,
all $33\,867$ labeled graphs on one through six vertices were screened in
double precision, covering $453\,098$ partial sums. The graph screen tested
the additive-one bound and the intermediate inequalities with tolerance
$10^{-6}$; all tests passed. It also included 32 larger structured graphs
and seven exact characteristic-polynomial checks. These numerical screens
are diagnostics, not certificates of a universal statement.

A supplementary search addresses the original question on a fixed
five-vertex universe, allowing unused vertices and faces of every dimension.
Recursion through faces in increasing cardinality, admitting a face only
after its boundary, gives $7\,579$ nonempty complexes. A separate
enumeration of maximal-face antichains gives the same set. Full vertex
permutation orbits reduce this set to 208 isomorphism classes. Thus no
assumption of flagness or of having exactly five used vertices enters
this finite search.

For each representative, integral matrices were reconstructed directly
from the face definitions, and all roots of their characteristic polynomials
were isolated with rational endpoints and multiplicities. The computation
uses SymPy~1.14.0 \cite{Meurer2017}; its real-root isolation routine implements
a continued-fraction method of the type described by Akritas and
Strzebo\'nski \cite{Akritas2005}. If the descending eigenvalues lie in
$[a_i,b_i]$ for $L$ and $[c_i,d_i]$ for $D$, the certificate checks
\[
 S_k(L)-S_k(D)\geq\sum_{i=1}^k(a_i-d_i)>0.
\]
All $3\,151$ representative partial sums, representing $113\,716$ labeled
partial sums, satisfy this strict inequality. The smallest recorded
rational lower bound is $1583407980/1583407981$. A separately implemented
verifier reconstructed the orbit coverage and characteristic polynomials
and checked the intervals by exact rational Sturm counts. This is a
strictly positive finite observation; the stated lower bound is less
than one and does not certify the stronger additive-one assertion in
higher dimensions.

Rational separation of algebraic spectral quantities also appears in
Zeng, Liu and Ratnavelu's exact Schatten-norm counterexample
\cite[Section 2, equations (8)--(18)]{Zeng2026}. Their comparison uses
rational bounds on radicals; the present finite search uses rational
root intervals for integral characteristic polynomials. In both settings
the sign decision is exact, whereas decimal approximations serve only
as diagnostics.

The accompanying computational files record the face masks, integral
matrices, characteristic polynomials, rational intervals and partial-sum
lower bounds, together with the enumeration and checking programs.
They provide finite evidence for the remaining problem. Extending
Theorem~\ref{thm:main} to complexes with higher-dimensional faces requires
an argument beyond the graph incidence blocks in \eqref{eq:blocks}; no
such extension or general obstruction is asserted here.

\paragraph{Computational and AI assistance.}
OpenAI Codex assisted with proof development, Python implementation,
literature searches and manuscript drafting. A separate model context
checked the argument and the supplementary certificates. This was not
external peer review or proof-assistant formalization. The numerical
screens used NumPy~2.2.6; exact polynomial and rational computations used
SymPy~1.14.0.


\raggedright
\bibliographystyle{plain}
\bibliography{references}

\end{document}
